\documentclass[11pt]{article}
\usepackage{fullpage,graphicx,hyperref,amsmath,amssymb}

\newcommand{\fig}[2]{\begin{figure}[h]\begin{center}%
  \includegraphics[scale=#2]{#1.pdf}\end{center}%
  \end{figure}}
\newcommand{\ceil}[1]{\left\lceil #1\right\rceil}

\begin{document}

\begin{center}\Large\bf 
CS 473 (Spring 2025)\\
Homework 9 (due Apr 24 Thu 10am)
\end{center}

\medskip\noindent
{\bf Instructions:} As in previous homeworks.

\begin{description}

\bigskip
\item[Problem 9.1:]
Consider the following problem {\sc Colorful-Path}:
\begin{quote}
\emph{Input}: a directed acyclic graph (dag) $G=(V,E)$ with $n$ vertices and $m$ edges, and
vertices $s,t\in V$,
and a color $c(e)\in \{1,\ldots,M\}$ for each edge $e\in E$, and a number~$\ell$.\\[2pt]
\emph{Output}: yes iff there exists a path from $s$ to $t$ that encounters at least $\ell$
\emph{distinct} colors.
\end{quote}

Prove that {\sc Colorful-Path} is NP-hard, by giving a polynomial-time reduction from
{\sc 3-SAT} to {\sc Colorful-Path}.  Remember to describe precisely your construction
of an input to {\sc Colorful-Path} from any given input to {\sc 3-SAT} (a 3CNF formula -- note that we are not given a satisfying assignment itself\@!), and carefully prove correctness of your reduction.

(Hint: for your construction of the dag $G$, try something like the picture below\ldots\ \ Describe how to assign colors to the edges along these paths.)


\item[Problem 9.2:]
Consider the following problem {\sc Covering-Square-by-Triangles}:
\begin{quote}
\emph{Input}: a set $T$ of $n$ triangles in 2D, a square $s$, and an integer $k$.\\[2pt]
\emph{Output}: yes iff there exists a subset $Q\subseteq T$ of $k$ triangles such that the union of $Q$ covers the entire boundary of $s$.
\end{quote}

(Note: we don't require the union of $Q$ to cover the interior of $s$.  Also, the variant of the problem of covering just one side of $s$, rather than the entire boundary of $s$, is easier and can  actually be solved in polynomial time by a greedy algorithm.)

Prove that {\sc Covering-Square-by-Triangles} is NP-hard.

You may use the fact that {\sc Vertex-Cover} is NP-hard even when restricted to graphs with maximum degree 3.  You may also use the fact that any graph with maximum degree $d$ is $(d+1)$-edge-colorable, i.e., we can assign a color from $\{1,\ldots,d+1\}$ to each edge, so that no two adjacent edges have the same color; such a coloring can be found in polynomial time.

As always, describe the construction of your reduction precisely, and carefully prove correctness of your reduction.

(Hint: map each edge of the input graph to a tiny interval along the boundary of $s$\ldots)


\bigskip
\item[Problem 9.3:]
As we all know, an \emph{interval} is a set $[a,b]$ containing all real numbers between $a$ and $b$.
Define a \emph{circular interval} to be either an interval $[a,b]$ (if $a\le b$),
or the union of two unbounded intervals $[a,\infty)\cup (-\infty,b]$ (if $a>b$).
(In other words, we allow the real line to ``wrap around''.)

Consider the following problem {\sc Circular-Interval-Selection}:
\begin{quote}
\emph{Input}: $N$ circular intervals $I_1,I_2,\ldots,I_N$ with positive integer weights $w_1,\ldots,w_N$.\\[2pt]
\emph{Output}: a subset $S\subseteq \{1,\ldots,N\}$ maximizing total weight $\sum_{i\in S}w_i$
such that for every $i,j\in S\ (i\neq j)$, the circular intervals $I_i$ and $I_j$ are disjoint.
\end{quote}

It is possible to solve the {\sc Circular-Interval-Selection} problem in $O(N^2)$ time by
dynamic programming (you don't have to give such an algorithm).  In contrast, without circularity, the problem can
be solved in $O(N\log N)$ time.  In this question, you will provide evidence on why the circular
version of the problem is more difficult.

To this end, consider the following known graph problem called {\sc Min-Weight-Triangle}:
\begin{quote}
\emph{Input}: a directed graph $G=(V,E)$ with $n$ vertices, where each edge $(u,v)$ has a positive  integer weight $w(u,v)$.\\[2pt]
\emph{Output}: a cycle of length 3 in $G$ with the minimum total weight, i.e., $v_i,v_j,v_k\in V$
with $(v_i,v_j),(v_j,v_k),(v_k,v_i)\in E$ minimizing $w(v_i,v_j)+w(v_j,v_k)+w(v_k,v_i)$.
\end{quote}
It has been hypothesized that no algorithm can solve {\sc Min-Weight-Triangle}
in $O(n^{2.999})$ time (in other words, the obvious cubic-time, brute-force algorithm is close to be the best possible).

Prove that under this hypothesis, no algorithm can solve {\sc Circular-Interval-Selection}
in $O(N^{1.499})$ time.

(Hint: map each edge $(v_i,v_j)$ of the input graph to three circular intervals $[i,M+j-1]$, $[M+i,2M+j-1]$, and $[2M+i,\infty)\cup (-\infty,j-1]$ for some $M$.  But how to assign weights to these circular intervals?)

\end{description}
\end{document}